\section{Introduction}
Federated learning distributes model training while retaining examples at
participating clients~\cite{mcmahan2017}. Its aggregation interface also gives
compromised participants a direct mechanism for modifying the global model.
In decentralized settings, the aggregation policy must additionally determine
who can participate, how updates are validated, and what happens after harmful
updates have been accepted.

Enrollment, robust aggregation, private evaluation, and recovery have distinct
trust boundaries. An enrolled participant can be malicious, heterogeneous data
can violate aggregation assumptions, and a ledger cannot establish the correctness
of private evaluations. Restoring model parameters also leaves earlier approval
transactions intact.

We study a two-layer architecture that makes these boundaries explicit. Layer~1
uses enrolled, staked validators to assess candidate updates on private local
validation data. Layer~2 monitors model quality and records checkpoint restoration
and attributed penalties. The ledger remains append-only throughout recovery.
The term \emph{accountable} denotes recorded responsibility and enforceable stake
accounting under the stated trust assumptions; it does not imply that a negative
quality trend proves malicious intent.

The intended contribution is an executable system specification with observable
failure boundaries: exact committee and quorum rules, disjoint validation and
test data, recovery that preserves blockchain history, and separate measurements
of malicious-update acceptance, downstream attack success, and wrongful penalties.
Component ablations and matched aggregation baselines evaluate the integrated
mechanism. EMH-DNN routing is a secondary case study. Robustness claims require
completed, matched-seed experiments.
